\documentclass[UTF8]{ctexart}
\usepackage{xeCJK}
\usepackage{geometry}
\geometry{left=2cm,right=2cm,top=2.5cm,bottom=2.5cm}
\setlength{\headheight}{12.7pt}

\usepackage{titlesec}
\titleformat{\section}
  {\normalfont\large\bfseries}
  {\thesection}
  {1em}
  {}
\titlespacing*{\section}
  {0pt}
  {3.5ex plus 1ex minus .2ex}
  {2.3ex plus .2ex}

\usepackage{amsmath}
\usepackage{amssymb}
\usepackage{amsthm}
\usepackage{enumitem}
\setlist[enumerate]{itemsep=0pt,topsep=0pt,parsep=0pt,leftmargin=0pt,itemindent=2.5em}
\setlist[itemize]{itemsep=0pt,topsep=0pt,parsep=0pt,leftmargin=0pt,itemindent=2.5em}
\usepackage{fancyhdr}
\pagestyle{fancy}
\fancyhf{}
\usepackage{graphicx}
\usepackage{hyperref}
\usepackage{indentfirst}
\usepackage{lmodern}

\begin{document}
\setlength{\abovedisplayskip}{1pt}
\setlength{\belowdisplayskip}{1pt}

\section{序型}

\hypertarget{sec:定义1.1}{}
\noindent\textbf{定义1.1 (序型)}

给定\href{../01 序关系/04 良序关系#nameddest=sec:定义1.1}{良序结构} $(a,\leqslant_a)$
和\href{01 序数#nameddest=sec:定义1.1}{序数} $\alpha$. 如果
\[
    (a,<_a)\cong(\alpha,\in),
\]
就称$\alpha$是$(a,\leqslant_a)$的\textbf{序型}, 记作$\mathrm{Type}(a,\leqslant_a)$.
在不产生歧义的情况下简记为$\mathrm{Type}(a)$.

显然良序结构的序型是唯一的.

\vspace{10pt}

\hypertarget{sec:定理1.1}{}
\noindent\textbf{定理1.1 (序型定理)}

任意的良序结构都有序型.

\noindent{\kaishu 证明.}

对任意的良序结构$(a,\leqslant_a)$, 取$a$中前段有序型的那些元素
\[
    b=\{x\in a\mid(a_x,\leqslant_a)\text{\,有序型\,}\},
\]
然后收集这些前段的序型
\[
    \alpha=\{\mathrm{Type}(a_x,\leqslant_a)\mid x\in b\}.
\]

首先, 如果$\gamma\in\beta\in\alpha$, 那么存在$x\in b$使得
\[
    \beta=\mathrm{Type}(a_x,\leqslant_a)\quad\implies\quad
    \exists f:(\beta,\in)\cong(a_x,<_a),
\]
进一步, $f\!\restriction_\gamma:(\gamma,\in)\cong(a_{f(\gamma)},<_a)$,
从而$\gamma\in\alpha$. 即$\alpha$是传递集. 从而$\alpha$是序数.

\vspace{3pt}
接下来, 根据良序集基本定理, 有以下的情况:

\begin{enumerate}
    \item 如果$a\cong\alpha$, 则$\alpha$是$(a,\leqslant_a)$的序型;
    \item 如果存在$x\in a$使得$a_x\cong\alpha$, 那么
    $\alpha\in\alpha$, 矛盾;
    \item 如果存在$\beta\in\alpha$使得$a\cong\beta$,
    那么存在$x\in a$使得$a_x\cong\beta\cong a$, 矛盾. \hfill $\qedsymbol$
\end{enumerate}

\vspace{10pt}

\hypertarget{sec:定理1.2}{}
\noindent\textbf{定理1.2}

任给良序结构$(a,\leqslant_a)$和$b\subset a$,
有$\mathrm{Type}(b)\leqslant\mathrm{Type}(a)$.

\noindent{\kaishu 证明.}

分别记$a,b$的序型为$\alpha,\beta$, 有序同构$f:\alpha\cong a$和$g:\beta\cong b$.

对$\gamma<\beta$归纳证明$\gamma\leqslant f^{-1}(g(\gamma))$.
任取$\gamma<\beta$, 假设对任意的$\delta<\gamma$
有$\delta\leqslant f^{-1}(g(\delta))$, 所以
\[
    g(\delta)<g(\gamma)\quad\implies\quad f^{-1}(g(\delta))<f^{-1}(g(\gamma))
    \quad\implies\quad \delta<f^{-1}(g(\gamma)),
\]
从而$\gamma\leqslant f^{-1}(g(\gamma))$.

接下来假设$\alpha<\beta$则有$\alpha\leqslant f^{-1}(g(\alpha))$,
但是$\mathrm{Dom}\,f=\alpha$, 矛盾. 所以$\beta\leqslant\alpha$. \hfill $\qedsymbol$

\vspace{10pt}

\hypertarget{sec:定理1.3}{}
\noindent\textbf{定理1.3}

任给某个良序结构的一族前段$S$, 有
$\displaystyle\mathrm{Type}\left(\bigcup S\right)=\sup_{s\in S}\mathrm{Type}(s)$.

\noindent{\kaishu 证明.}

取序同构$\displaystyle f:\bigcup S\cong\mathrm{Type}\left(\bigcup S\right)$, 则
\begin{equation*}
    \mathrm{Type}\left(\bigcup S\right)=f\left[\,\bigcup S\,\right]=
    \bigcup_{s\in S}f[s]=\sup_{s\in S}\mathrm{Type}(s).
\tag*{\qedsymbol}\end{equation*}

\end{document}
